\subsection{The normal extension theorem, including nonexcellent bases}

\begin{theorem}[C-RD6: the normal case]\label{thm:C-RD6}
\status{Proved by the checked normal-base extension theorem.}
Let $S$ be a locally noetherian normal scheme, let $U\subset S$ be a dense
open, and suppose every point of $S\setminus U$ has residue
characteristic $p$. The functor $F$ of Part C is an equivalence.
Neither excellence nor finiteness of normalizations of modifications of
$S$ is needed. In particular the normal case of C.2 holds without
hypotheses a) or b), apart from the indicated density.
\end{theorem}
\begin{proof}
Work on one connected component, which is integral. For a triple
$(A_U,M,\varphi)$, \cite[Corollary~6.12]{KrishnamoorthyPal2021RankTwo}
produces an abelian scheme $A/S$ extending $A_U$. This theorem has
exactly the hypotheses displayed here: locally noetherian and normal,
with the missing locus in characteristic $p$. It is not a purity
statement requiring a regular or unramified base.

A normal noetherian integral scheme is a Krull scheme. Apply
\cite[Theorem~6.3]{KrishnamoorthyPal2021RankTwo} to the generic-fibre
isomorphism $A[p^\infty]\simeq M$ induced by $\varphi$. That theorem
states the full faithfulness of restriction of Barsotti--Tate groups
from an integral Krull scheme to its generic point. Apply it to the
isomorphism and its inverse; their composites are identities by
faithfulness. This gives the required isomorphism over $S$. Its
restriction to $U$ is the original one, again by faithfulness. Thus the
entire triple, including its specified identification, is realized.
For full faithfulness, a morphism of triples gives a homomorphism of
BT groups over $S$ which is algebraic over the dense open $U$. Apply
\cite[Proposition~1.2, p.64]{Grothendieck1966Hom}: for every associated
point of the reduced scheme $S$, the required algebraicity hypothesis
holds at that generic point itself. The proposition produces the unique
abelian-scheme homomorphism inducing the prescribed BT map. Its
restriction to $U$ is the original homomorphism by faithfulness.
All constructions agree on overlaps by uniqueness and hence glue.

For clarity, the nonexcellent point is explicitly addressed in the
source's proof of Corollary~6.12. It takes the closure of the graph of a
map to a fine polarized moduli scheme, then passes to its normalization,
which need not be finite or noetherian. The normalization is an integral
Krull scheme, so Theorem~6.3 applies. The proof uses the Mori--Nagata
statements cited there as EGA~$0_{\mathrm{IV}}$, 23.2.6--23.2.8 to
control the reduced fibres, then Corollary~6.10 and Proposition~6.11 to
show that the graph has zero-dimensional fibres. This is precisely the
normalization issue that cannot be omitted by imposing an unstated
excellence assumption.
\end{proof}

\begin{remark}[An additional gap in the inherited spreading argument]
The proof in the supplied reading of 100:15--17 invokes EGA~IV, 8.8.2
to spread both an isomorphism of abelian schemes and an isomorphism of
$p$-divisible groups from a local scheme to a neighbourhood. The first
invocation is justified by finite presentation. The second is an
infinite compatible system and does not follow from that theorem.
One needs a separate finite-level determination or an algebraicity-locus
argument. Consequently that inherited argument is not used to remove
normality in Theorem~\ref{thm:C-RD6}.
\end{remark}
